% !TEX TS-program = XeLaTeX
% use the following command:
% all document files must be coded in UTF-8
\documentclass[portuguese]{textolivre}
% build HTML with: make4ht -e build.lua -c textolivre.cfg -x -u article "fn-in,svg,pic-align"

\journalname{Texto Livre}
\thevolume{19}
%\thenumber{1} % old template
\theyear{2026}
\receiveddate{\DTMdisplaydate{2026}{4}{8}{-1}} % YYYY MM DD
\accepteddate{\DTMdisplaydate{2026}{5}{15}{-1}}
\publisheddate{\DTMdisplaydate{2026}{5}{15}{-1}}
\corrauthor{Vânia Soares Barbosa}
\articledoi{10.1590/1983-3652.2026.65593}
%\articleid{NNNN} % if the article ID is not the last 5 numbers of its DOI, provide it using \articleid{} commmand 
% list of available sesscions in the journal: articles, dossier, reports, essays, reviews, interviews, editorial
\articlesessionname{dossier}
\runningauthor{Barbosa} 
%\editorname{Leonardo Araújo} % old template
\sectioneditorname{Zaira Bomfante dos Santos~\orcid{000-0002-6162-8489}}
\layouteditorname{Leonardo Araújo~\orcid{0000-0003-3884-2177}}

\title{Resenha de \textit{Territórios do Letramento}}
\othertitle{Review of \textit{Territórios do Letramento}}

% if there is a third language title, add here:
%\othertitle{Artikelvorlage zur Einreichung beim Texto Livre Journal}

\author[1]{Vânia Soares Barbosa~\orcid{0000-0002-5816-1064}\thanks{Email: \href{mailto:vaniasb@ufpi.edu.br}{vaniasb@ufpi.edu.br}}}
\affil[1]{Universidade Federal do Piauí, Programa de Pós-Graduação em Letras, Teresina, PI, Brasil.}


\addbibresource{article.bib}
% use biber instead of bibtex
% $ biber article

% used to create dummy text for the template file
\definecolor{dark-gray}{gray}{0.35} % color used to display dummy texts
\usepackage{lipsum}
\SetLipsumParListSurrounders{\colorlet{oldcolor}{.}\color{dark-gray}}{\color{oldcolor}}

% used here only to provide the XeLaTeX and BibTeX logos
\usepackage{hologo}

% if you use multirows in a table, include the multirow package
\usepackage{multirow}

% provides sidewaysfigure environment
\usepackage{rotating}

\usepackage{endnotes} % Inclui o pacote endnotes

% CUSTOM EPIGRAPH - BEGIN 
%%% https://tex.stackexchange.com/questions/193178/specific-epigraph-style
\usepackage{epigraph}
\renewcommand\textflush{flushright}
\makeatletter
\newlength\epitextskip
\pretocmd{\@epitext}{\em}{}{}
\apptocmd{\@epitext}{\em}{}{}
\patchcmd{\epigraph}{\@epitext{#1}\\}{\@epitext{#1}\\[\epitextskip]}{}{}
\makeatother
\setlength\epigraphrule{0pt}
\setlength\epitextskip{0.5ex}
\setlength\epigraphwidth{.7\textwidth}
% CUSTOM EPIGRAPH - END

% to use IPA symbols in unicode add
%\usepackage{fontspec}
%\newfontfamily\ipafont{CMU Serif}
%\newcommand{\ipa}[1]{{\ipafont #1}}
% and in the text you may use the \ipa{...} command passing the symbols in unicode

% LANGUAGE - BEGIN
% ARABIC
% for languages that use special fonts, you must provide the typeface that will be used
% \setotherlanguage{arabic}
% \newfontfamily\arabicfont[Script=Arabic]{Amiri}
% \newfontfamily\arabicfontsf[Script=Arabic]{Amiri}
% \newfontfamily\arabicfonttt[Script=Arabic]{Amiri}
%
% in the article, to add arabic text use: \textlang{arabic}{ ... }
%
% RUSSIAN
% for russian text we also need to define fonts with support for Cyrillic script
% \usepackage{fontspec}
% \setotherlanguage{russian}
% \newfontfamily\cyrillicfont{Times New Roman}
% \newfontfamily\cyrillicfontsf{Times New Roman}[Script=Cyrillic]
% \newfontfamily\cyrillicfonttt{Times New Roman}[Script=Cyrillic]
%
% in the text use \begin{russian} ... \end{russian}
% LANGUAGE - END

% EMOJIS - BEGIN
% to use emoticons in your manuscript
% https://stackoverflow.com/questions/190145/how-to-insert-emoticons-in-latex/57076064
% using font Symbola, which has full support
% the font may be downloaded at:
% https://dn-works.com/ufas/
% add to preamble:
% \newfontfamily\Symbola{Symbola}
% in the text use:
% {\Symbola }
% EMOJIS - END

% LABEL REFERENCE TO DESCRIPTIVE LIST - BEGIN
% reference itens in a descriptive list using their labels instead of numbers
% insert the code below in the preambule:
%\makeatletter
%\let\orgdescriptionlabel\descriptionlabel
%\renewcommand*{\descriptionlabel}[1]{%
%  \let\orglabel\label
%  \let\label\@gobble
%  \phantomsection
%  \edef\@currentlabel{#1\unskip}%
%  \let\label\orglabel
%  \orgdescriptionlabel{#1}%
%}
%\makeatother
%
% in your document, use as illustraded here:
%\begin{description}
%  \item[first\label{itm1}] this is only an example;
%  % ...  add more items
%\end{description}
% LABEL REFERENCE TO DESCRIPTIVE LIST - END


% add line numbers for submission
%\usepackage{lineno}
%\linenumbers

\begin{document}
\maketitle

\begin{figure}[htbp]
\centering
\begin{minipage}{.5\textwidth}
\includegraphics[width=\textwidth]{Fig1.png}
 \caption*{\fullcite{brambila2024territorios}.}
 %\caption{\fullcite{nascimento_formacao_2019}.}
 \label{fig01}
\end{minipage}
\end{figure}


O ano era 1994. Na cidade estadunidense de New London, um grupo de dez educadores dos Estados Unidos, Inglaterra e Austrália discutiam o futuro das prática pedagógicas frente aos desafios das tecnologias digitais de comunicação e informação. No Brasil, fomentavam as discussões sobre alfabetização e letramento, iniciadas desde os tempos de Paulo Freire e fortalecidas com as ideias de Magda Soares, entre outros educadores e pesquisadores nacionais. Estes cenários, que parecem separados geograficamente, na verdade, trazem em comum a base sólida, mas não imutável, da busca por uma educação inclusiva e um lugar no qual apenas saber escrever e ler, como codificar e decodificar signos linguísticos, já não atendem mais às demandas da sociedade contemporânea. É preciso formar indivíduos sociais autonômos e aptos a manusearem a pluralidade dos textos que circulam nessa sociedade, que, como o livro organizado por Guilherme Brambila vai mostrar, constituem-se em “Territórios do Letramento” nos quais aqueles indivíduos são agentes atuantes na e por meio da linguagem. O convite a um passeio por esses territórios é o que o leitor do livro irá encontrar desde a sua capa, passando por todos os elementos que o constitui.

Em uma leitura multimodal que integra todos os modos e recursos semióticos no texto, como uma porta de entrada, a capa traz a imagem de duas figuras femininas – uma criança e uma mulher – envolvidas no ato da leitura de um livro. Nesse ato, a produção de significados é visualmente representada por borboletas e outros elementos que compõem o cenário mágico que as envolve, com um destaque às cores em tons de amarelo e vinho, podendo trazer a ideia das diferentes terras sendo iluminadas pelo sol. As possibilidades de leitura dessa capa não são únicas, mas, de certa forma, materializam a ideia

\begin{quote}
     [d]a criança que ainda não se alfabetizou, mas já folheia livros, finge lê-los, brinca de escrever, ouve histórias que lhe são lidas, está rodeada de material escrito e percebe seu uso e função, essa criança é ainda "analfabeta", porque não aprendeu a ler e a escrever, mas já penetrou no mundo do \textbf{letramento}, já é, de certa forma, \textbf{letrada} \cite[p. 24, grifos nossos]{soares2025letramento}.
\end{quote}

Ainda que epistologicamente a língua portuguesa já favoreça a distinção entre alfabetização e letramento, e para além de um certo protagonismo da educadora Magda Soares, cujas ideias preconizam dois dos seis capítulos do livro, a escolha da capa certamente representa o ponto de partida para o que o leitor irá encontrar em todo o escopo da obra, cujos elementos internos se constituem na Apresentação, feita pelo próprio organizador, e em seis capítulos de autoria de pesquisadores de diferentes regiões do grande território brasileiro, “detentores de uma produção consistente e altamente contribuitiva aos estudos do letramento” \cite[p. 9]{brambila2024apresentacao}.  

Na Apresentação, \textcite{brambila2024apresentacao} explica os objetivos da obra, traz um resumo dos territórios dos quais os capítulos irão tratar, e deixa claro que estes não são terrítórios únicos e que, portanto, a pergunta “quais são os territórios do letramento?” encontra sim o fortalecimento nas ideias discutidas no livro, mas não uma resposta fechada e definitiva. Tal afirmação não deixa de provocar no leitor a inquietação e o desejo por descobrir mais sobre esses territórios nos quais, manifestadas nos usos sociais da língua, habitam práticas sociais diversas. Essa busca, que se inicia nos capítulos que seguem, percorerá os caminhos dos estudos textuais, em suas diferentes perspectivas – discursivas, semióticas e tradutológicas – e relações com as tecnologias digitais e de comunicação. O resultado dessa travessia, como o organizador manifesta seu desejo, mas não podendo ser diferente, verberará na formação de professores de línguas e, consequentemente, na formacão do cidadão contemporâneo. 

A metáfora do livro como uma travessia, que adotamos nesta resenha, traz, de certa forma, um tom de ancestralidade às concepções de texto e letramento de Luiz Antônio Marcuschi e Magda Soares, respectivamente, além de Ingedore Koch, Roxane Rojo, Ângela Kleiman, Gunther Kress e Theo van Leeuwen, entre outros. 

O conjunto dessas “vozes discursivas” permeia os capítulos da obra, inciando com o capítulo \textit{Texto, sentido e letramento}, de autoria do próprio organizador,  no qual o autor \cite{brambila2024texto} propõe o não enclausuramento do texto, deixando-o aberto para interações entre indivíduos sociais e, juntamente com o campo dos letramentos, pensá-los como “um elo produtivo à compreensão da linguagem, das práticas sociais e do complexo processo de tornar-se letrado contemporaneamente” \cite[p. 13]{brambila2024texto}. Para tanto, Brambrilha organiza o capítulo em três seções: na primera, aquele não enclausuramento e sua relação próxima com as interações humanas é proposto a partir das ideias prioneiras de práticas de letramento de Ângela Kleiman e da construção de sentidos no texto de Ingedore Koch. Essa relação intrísica entre a construcao de sentidos no texto e as práticas de letramento é construída a partir de bases teóricas sólidas que, juntando-se ao tom autoral proporcionado por exemplos retirados da tese do próprio autor, constituindo-se em um ponto forte da seção. Dando sequência, na segunda seção o leitor é convidado a revisitar, ou reaprender, ainda que, como o autor diz, suscitamente, as contribuições de Luiz Antônio Marcushi para a Linguística Textual, mantendo-se a perspectiva do elo entre texto e letramento e suas implicações para as práticas sociais e pedagógicas.  A travessia entre as seções finaliza como uma extensão da Apresentação, agora com o suporte do arcabouço teórico discursivo apresentado nas seções anteriores e marcando as esferas iniciais dos territórios do letramento.

No capítulo seguinte  - \textit{Letramentos acadêmicos: enfoques sobre a escrita científica}, as autoras Adriana Fisher, Letícia Lungen, Mariana Aparecida Vicentini e Sandra Pottmeier, nos convidam a pensarmos os letramentos no ambiente acadêmico, olhando para a escrita acadêmica como um espaço de vozes discursivas em diferentes areas do conhecimento, cada uma com suas particularidades. A travessia no território do letramento acadêmico é feita a partir de um estudo qualitativo que investigou a autocitação em artigos científicos em duas distintas áreas – Psicologia e Biologia – buscando compreender suas funções e sentidos nas produções desse gênero discursivo. Para além da discussão sólida dos resultados obtidos, destaca-se desse capítulo o percurso metodológico apresentado pelas autoras que, a despeito de seus ‘percalços’, como as próprias se referem, pode servir de modelo para outros estudos da mesma natureza.

Da academia para a infância, o capítulo seguinte traz o ineditismo no olhar para o brinquedo   como objeto de práticas semióticas multifacetadas que resifignificam o ato de brincar ao inseri-lo como prática social discursiva e, como tal, seu consumo e disseminação precisam ser devidamente discutidos e entendidos. Essa travessia é feita com o suporte dos estudos em multimodalidade, na Gramática do Design Visual, de Kress e van Leeuwen, e, em particular, do letramento do brinquedo, ideia teórica de Danielle Almeida que, juntamente com José Maria Sarinho Júnior, assinam o capítulo \textit{Letramento do brinquedo: deficiência, diferença e inclusão}. Neste capítulo, os autores olham para brinquedos artesanais inclusivos – bonecos fabricados artesanalmente que trazem em suas características a representação de pessoas com algum tipo de deficiência física – buscando entender “como são construídos significados interativos e composicionais sobre a deficiência, a diferença e a inclusão” nestes objetos. Os autores concluem mostrando como características sensoriais e afetivas podem  nos afetar e nos envolver em diferentes aspectos – sociais, históricos, culturais e ideológicos. 

Até aqui, a travessia nos capítulos trazem em comum o suporte de uma base teórica e metodológica ao, ainda que em diferentes proporções, apresentarem exemplos e/ou resultados de pesquisas desenvolvidas pelos próprios autores. Diferente do capítulo que se segue, no qual a escrita leve e ao mesmo tempo provocativa de Ana Elisa Ribeiro proporciona um passeio panorâmico no cenário brasileiro das tecnologias digitais de comunicação e informação, desde a chegada dos primeiros computadores nos anos 90, passando pelos territórios da Linguística e da Linguística Aplicada e dos documentos oficiais, chegando nas discussões que diferenciam alfabetização de letramento, já protagonizadas por mulheres desde 1980, especialmente por Magda Soares, homenageada no capítulo a partir do resgate de seu texto \textit{Novas práticas de leitura e escrita: letramento na cibercultura} \cite{soares2002cibercultura} e finalizando com o termo  “multiletramentos” preconizado no manifesto do New London Group (1996). Do seu título – \textit{Letramento Digital: definições impermanentes} – o capítulo chama a atenção para a incompletude do entendimento não apenas dos termos citados, ainda que sejam o ponto principal, mas, em especial, das práticas sociais mediadas pelas technologias digitais que ora são vistas como benéficas, ora como desafiadoras, mas sempre apresentando questões impermanentes e que acabam por revelar desigualdades sociais, políticas e econômicas.  

Em uma travessia semelhante, no capítulo \textit{Tradução, letramento e suas relações}, Patrick Rezende inicia fazendo um resgate, e reflexões, das definições de letramento e seus afastamentos/aproximações de alfetização, desde aqueles presentes em dicionários aos propostos pelas teóricas Mary Kato, Ingedore Kock, Ângela Kleiman e, assim como Ana Elisa Ribeiro, dá o merecido protagonismo à Magda Soares e seu livro \textit{Letramento: um tema em três gêneros} \cite{soares2025letramento}, cuja primeira edição foi publicada em 1998. Ao conjunto dessas vozes femininas o autor junta a voz de Paulo Freire que, já na década de setenta, chamava a atenção para algo que está além do saber ler e escrever, algo que já fazia com que a alfabetização fosse vista como “um ato político” \cite[p. 70]{freire1978alfabetizacao} e, como prática social que requer “a particição ativa e crítica dos sujeitos” \cite[p. 94]{rezende2024traducao}.  Essa participação é apresentada por Rezende como em uma viagem, levando o leitor aos territórios do letramento no campo da tradução no qual “o viajante, frequentemente, experimenta alguns estranhamentos, percebendo detalhes antes negligenciados na paissagem local” \cite[p. 97]{rezende2014ponte}, no ir e vir como em uma “ponte plástica” \cite[p. 97]{rezende2014ponte} que conecta, mas não limita as trocas interculturais e linguísticas na “construção de um mundo mais equitativo” \cite[p. 106]{rezende2014ponte}.

Trocas possivelmente é a definição mais próxima para o capítulo que, apenas fisicamente, encerra o livro aqui resenhado. Ao chegar nessa parte dos Territórios, Daniel Ferraz, autor deste último capítulo, nos convida para \textit{Uma breve conversa sobre linguagem, línguas e Linguística Aplicada}, com trocas intergeracionais ao homenagear sua avó que, segundo ele, “era letrada pela vida e possuía uma sabedoria insipiradora” \cite[p. 109]{feraz2024linguagem}.  À imagem dessa figura feminia sábia e ancestral, que abre o capítulo, Ferraz apresenta uma série de outras imagens e, a partir delas, propõe pausas em caixas com perguntas que objetivam refletir sobre os territórios dos letramentos visuais, das visões de língua, da educação linguística, da Linguística Aplicada, da Linguística Aplicada Crítica, e de outras questões que norteiam a sociedade contemporânea e suas complexidades. Como que reabrindo os territórios do letramento já visitados nos capítulos anteriores, Daniel Ferraz encerra o seu capítulo com a voz de Paulo Freire que exalta a necessidade do dialógo entre os homens para sua existencialidade – “[…] Se ao dizer suas palavras, ao chamar ao mundo, os homens o transformam, o dialógo impõe-se como o caminho pelo qual os homens encontram seu significado enquanto homens […]” \cite[p. 42]{freire1980educacao} – estendendo essa “necessidade existencial para todos os seres vivos deste planeta” \cite[p. 121]{feraz2024linguagem}. 

O livro \citetitle{brambila2024territorios} é finalizado com uma breve biografia do organizador e demais autoras e autores que assinam seus capítulos. Mas a travessia não se finda nessas páginas ou mesmos nos textos apresentados nas orelhas e na contracapa, pelo contrário, ela apenas fomenta a pergunta “quais são os territórios do letramento?”, como já mencionado no início desta resenha, e esta, por sua vez, é fometada, inquietantemente, pelas de Ana Elisa Ribeiro: “As perguntas de vinte anos atrás são as mesmas de hoje? Encontramos algumas repostas nesse intervalo? Quais? Que novas inquietações temos hoje?” \cite[p. 86]{ribeiro2024territorios}. A quem busca respostas prontas para aquela e essas perguntas, no livro não irá encontrar. Mas a diversidade de territórios, aqui discutidas com o suporte de um protagonismo feminino e ideias ancestrais de um tempo bem próximo, como as de Freire e Marcuschi, é certamente um convite a novos desbravamentos e outras tantas travessias no mundo das linguagens que permeiam todas as práticas sociais do presente século. 

\begin{spacing}{0.9}
\printbibliography\label{sec-bib}
\end{spacing}
% if the text is not in Portuguese, it might be necessary to use the code below instead to print the correct ABNT abbreviations [s.n.], [s.l.]
%\begin{portuguese}
%\printbibliography[title={Bibliography}]
%\end{portuguese}

\section*{\fontsize{10}{12}\selectfont Agradecimento}
A autora agradece à Editora, ao organizador do livro – Guilherme Brambila – e aos autores dos capítulos, em especial a Daniellle Almeida (in Memorian) e Ana Elisa Ribeiro, pelo envio do exemplar do livro “Territórios do Letramento”.

\begin{dataavailability}
\txtdataavailability{nodata} % options: dataavailable, dataonly, databody, datanotav, nodata
\end{dataavailability}

\end{document}

